% GENERATED FILE. Edit canonical lessons, then run npm run generate:books.
% canonical-source-hash: fnv1a64:05deb2a844bd5b60
% canonical-lessons: SA-C504-avedanam, SA-C504-niyuktih, SA-C504-pratibha, SA-C504-upajivati, SA-C504-pravinah

\chapter{Looking for Work}
\label{ch:sa-looking-for-work}
\hlchaptermodality{504}

\begin{chapteropening}
\textbf{By the end of this chapter:} \emph{I can apply for a job and say what I am good at.}

Complete the last lesson of chapter 504: I can apply for a job and say what I am good at.
\end{chapteropening}

\section[\texorpdfstring{\sk{आवेदनम्} (āvedanam)}{आवेदनम् (āvedanam)}]{\sk{आवेदनम्} (āvedanam) — an application}

\label{lesson:SA-C504-avedanam}



\begin{quote}
Before the new one: say the Sanskrit for an employer, then the Sanskrit for employs.
\end{quote}

\subsection*{You'll want to know: \sk{आवेदनम्}}
\textbf{\sk{आवेदनम्}} — \emph{āvedanam} — \textquotedblleft{}an application\textquotedblright{}.

\begin{quote}
\sk{अहं} \sk{कार्यालयं} \sk{प्रति} \sk{आवेदनं} \sk{प्रेषयामि।}
\end{quote}

\begin{quote}
\emph{ahaṃ kāryālayaṃ prati āvedanaṃ preṣayāmi.}
\end{quote}

\begin{quote}
I send an application to the office.
\end{quote}

A neuter noun.

\begin{cousinweb}
\textbf{\sk{आवेदनम्}} (\emph{āvedanam}) comes from \emph{ā} and the causative of the root √vid, \textquotedblleft{}to know\textquotedblright{}: a making known; MW gives \textquotedblleft{}announcing\textquotedblright{} and \textquotedblleft{}addressing respectfully\textquotedblright{}; \textquotedblleft{}an application for a post\textquotedblright{} is the modern sense (MW).
\end{cousinweb}

\begin{grammarlens}[title={What you've built}]
One more word for describing what you lived and felt.
\end{grammarlens}

\subsection*{Your turn}
\begin{itemize}
\raggedright
  \item \emph{Say it:} \emph{āvedanam}
  \item \emph{Say it:} \emph{āvedanam}, once more
  \item \emph{Say it:} \emph{ahaṃ kāryālayaṃ prati āvedanaṃ preṣayāmi.} — I send an application to the office.
  \item \emph{Recall:} say the Sanskrit for an employer, then the Sanskrit for employs, then say \emph{āvedanam} again
\end{itemize}

\begin{tcolorbox}[breakable,colback=teal!4,colframe=teal!35!black,title={Before you move on}]
What does \sk{आवेदनम्} mean? (\textquotedblleft{}An application\textquotedblright{}.) Say it once more.
\end{tcolorbox}
\section[\texorpdfstring{\sk{नियुक्तिः} (niyuktiḥ)}{नियुक्तिः (niyuktiḥ)}]{\sk{नियुक्तिः} (niyuktiḥ) — a post; an appointment}

\label{lesson:SA-C504-niyuktih}



\begin{quote}
Before the new one: say the Sanskrit for employs, then the Sanskrit for an application.
\end{quote}

\subsection*{You'll want to know: \sk{नियुक्तिः}}
\textbf{\sk{नियुक्तिः}} — \emph{niyuktiḥ} — \textquotedblleft{}a post; an appointment\textquotedblright{}.

\begin{quote}
\sk{मम} \sk{भ्रातुः} \sk{नूतना} \sk{नियुक्तिः} \sk{अस्ति।}
\end{quote}

\begin{quote}
\emph{mama bhrātuḥ nūtanā niyuktiḥ asti.}
\end{quote}

\begin{quote}
My brother has a new post.
\end{quote}

A feminine noun.

\begin{cousinweb}
\textbf{\sk{नियुक्तिः}} (\emph{niyuktiḥ}) is built on \emph{ni} and the root √yuj, \textquotedblleft{}to yoke, to appoint\textquotedblright{}, with the ending \emph{-ti}; MW gives \emph{niyukti} as an injunction, a charge, an office (MW).
\end{cousinweb}

\begin{grammarlens}[title={What you've built}]
One more word for describing what you lived and felt.
\end{grammarlens}

\subsection*{Your turn}
\begin{itemize}
\raggedright
  \item \emph{Say it:} \emph{niyuktiḥ}
  \item \emph{Say it:} \emph{niyuktiḥ}, once more
  \item \emph{Say it:} \emph{mama bhrātuḥ nūtanā niyuktiḥ asti.} — My brother has a new post.
  \item \emph{Recall:} say the Sanskrit for employs, then the Sanskrit for an application, then say \emph{niyuktiḥ} again
\end{itemize}

\begin{tcolorbox}[breakable,colback=teal!4,colframe=teal!35!black,title={Before you move on}]
What does \sk{नियुक्तिः} mean? (\textquotedblleft{}A post; an appointment\textquotedblright{}.) Say it once more.
\end{tcolorbox}
\section[\texorpdfstring{\sk{प्रतिभा} (pratibhā)}{प्रतिभा (pratibhā)}]{\sk{प्रतिभा} (pratibhā) — talent; a gift}

\label{lesson:SA-C504-pratibha}



\begin{quote}
Before the new one: say the Sanskrit for an application, then the Sanskrit for a post.
\end{quote}

\subsection*{You'll want to know: \sk{प्रतिभा}}
\textbf{\sk{प्रतिभा}} — \emph{pratibhā} — \textquotedblleft{}talent; a gift\textquotedblright{}.

\begin{quote}
\sk{सर्वे} \sk{तस्य} \sk{प्रतिभां} \sk{प्रशंसन्ति।}
\end{quote}

\begin{quote}
\emph{sarve tasya pratibhāṃ praśaṃsanti.}
\end{quote}

\begin{quote}
Everyone praises his talent.
\end{quote}

A feminine noun.

\begin{cousinweb}
\textbf{\sk{प्रतिभा}} (\emph{pratibhā}) is built on \emph{prati}, \textquotedblleft{}towards\textquotedblright{}, and the root √bhā, \textquotedblleft{}to shine, to appear\textquotedblright{}: what lights up in the mind; MW gives \emph{pratibhā} as light, then intelligence, presence of mind, genius (MW).
\end{cousinweb}

\begin{grammarlens}[title={What you've built}]
One more word for describing what you lived and felt.
\end{grammarlens}

\subsection*{Your turn}
\begin{itemize}
\raggedright
  \item \emph{Say it:} \emph{pratibhā}
  \item \emph{Say it:} \emph{pratibhā}, once more
  \item \emph{Say it:} \emph{sarve tasya pratibhāṃ praśaṃsanti.} — Everyone praises his talent.
  \item \emph{Recall:} say the Sanskrit for an application, then the Sanskrit for a post, then say \emph{pratibhā} again
\end{itemize}

\begin{tcolorbox}[breakable,colback=teal!4,colframe=teal!35!black,title={Before you move on}]
What does \sk{प्रतिभा} mean? (\textquotedblleft{}Talent; a gift\textquotedblright{}.) Say it once more.
\end{tcolorbox}
\section[\texorpdfstring{\sk{उपजीवति} (upajīvati)}{उपजीवति (upajīvati)}]{\sk{उपजीवति} (upajīvati) — lives by; earns a living from}

\label{lesson:SA-C504-upajivati}



\begin{quote}
Before the new one: say the Sanskrit for a post, then the Sanskrit for talent.
\end{quote}

\subsection*{You'll want to know: \sk{उपजीवति}}
\textbf{\sk{उपजीवति}} — \emph{upajīvati} — \textquotedblleft{}lives by; earns a living from\textquotedblright{}.

\begin{quote}
\sk{सः} \sk{वेतनेन} \sk{उपजीवति।}
\end{quote}

\begin{quote}
\emph{saḥ vetanena upajīvati.}
\end{quote}

\begin{quote}
He lives on his wages.
\end{quote}

Past (he …): \sk{सः} \sk{वेतनेन} \sk{उपजीवितवान्}, \emph{saḥ vetanena upajīvitavān}.

\begin{cousinweb}
\textbf{\sk{उपजीवति}} (\emph{upajīvati}) is built on \emph{upa}, \textquotedblleft{}near, under\textquotedblright{}, and the root √jīv, \textquotedblleft{}to live\textquotedblright{}, which also gives \emph{jīvati}; MW gives \emph{upajīv} as to live upon, to subsist by, to make use of (MW).
\end{cousinweb}

\begin{grammarlens}[title={What you've built}]
One more word for describing what you lived and felt.
\end{grammarlens}

\subsection*{Your turn}
\begin{itemize}
\raggedright
  \item \emph{Say it:} \emph{upajīvati}
  \item \emph{Say it:} \emph{upajīvati}, once more
  \item \emph{Say it:} \emph{saḥ vetanena upajīvati.} — He lives on his wages.
  \item \emph{Recall:} say the Sanskrit for a post, then the Sanskrit for talent, then say \emph{upajīvati} again
\end{itemize}

\begin{tcolorbox}[breakable,colback=teal!4,colframe=teal!35!black,title={Before you move on}]
What does \sk{उपजीवति} mean? (\textquotedblleft{}Lives by; earns a living from\textquotedblright{}.) Say it once more.
\end{tcolorbox}
\section[\texorpdfstring{\sk{प्रवीणः} (pravīṇaḥ)}{प्रवीणः (pravīṇaḥ)}]{\sk{प्रवीणः} (pravīṇaḥ) — skilled; versed}

\label{lesson:SA-C504-pravinah}



\begin{quote}
Before the new one: say the Sanskrit for talent, then the Sanskrit for lives by.
\end{quote}

\subsection*{You'll want to know: \sk{प्रवीणः}}
\textbf{\sk{प्रवीणः}} — \emph{pravīṇaḥ} — \textquotedblleft{}skilled; versed\textquotedblright{}.

\begin{quote}
\sk{सा} \sk{संगणके} \sk{प्रवीणा} \sk{अस्ति।}
\end{quote}

\begin{quote}
\emph{sā saṅgaṇake pravīṇā asti.}
\end{quote}

\begin{quote}
She is skilled with computers.
\end{quote}

\begin{cousinweb}
\textbf{\sk{प्रवीणः}} (\emph{pravīṇaḥ}) is built on \emph{pra} and \emph{vīṇā}, \textquotedblleft{}the lute\textquotedblright{}, as MW gives it; MW gives \emph{pravīṇa} as skilful, clever, versed in a thing (MW).
\end{cousinweb}

\begin{grammarlens}[title={What you've built}]
One more word for describing what you lived and felt.
\end{grammarlens}

\subsection*{Your turn}
\begin{itemize}
\raggedright
  \item \emph{Say it:} \emph{pravīṇaḥ}
  \item \emph{Say it:} \emph{pravīṇaḥ}, once more
  \item \emph{Say it:} \emph{sā saṅgaṇake pravīṇā asti.} — She is skilled with computers.
  \item \emph{Recall:} say the Sanskrit for talent, then the Sanskrit for lives by, then say \emph{pravīṇaḥ} again
\end{itemize}

\begin{tcolorbox}[breakable,colback=teal!4,colframe=teal!35!black,title={Before you move on}]
What does \sk{प्रवीणः} mean? (\textquotedblleft{}Skilled; versed\textquotedblright{}.) Say it once more.
\end{tcolorbox}
